\begin{abstract}
This note gives a complete account of quantitative parameter mapping with the multi-pathway
multi-echo (MPME) sequence of Cheng et al.~\cite{cheng2019}. We derive the steady-state
forward model from the configuration-state description and prove that, for each scan, the
signal depends on the flip angle and $\Tone$ only through two invariants
(\cref{thm:uv}); one scan therefore cannot separate $\Mzero$, $\Bone$ and $\Tone$, and for
small flip angles two scans cannot either, to leading order. The published protocol escapes
this through its near-$360^\circ$ second pulse. We implement the published analytic
inverse and verify it, analyse the Fisher information and Cram\'er--Rao bounds (CRLB) of the
model, identify its exact global aliases, and show that a joint maximum-likelihood (ML) fit
attains the bound where the bound applies, while the analytic inverse is 3--7 times above
it. For magnitude data we reduce the alias conditions to a square system and show that an exact
alias exists inside the physiological $\Bone$ range and explain why a magnitude fit needs the relative phase of the two FID images (one
bit) while complex ML does not; with that bit, magnitude least squares performs like complex
ML. Exact Jacobians obtained by implicit differentiation of the steady state make both fits
10--34 times faster, so that a $256\times256$ slice is fitted in seconds to minutes on a laptop
CPU. A two-dimensional digital phantom with simulated acquisitions demonstrates all
estimators, including the published low-resolution second scan, and validates per-voxel
uncertainty maps; the BrainWeb phantom quantifies how partial volume with CSF biases every
single-compartment estimator while leaving the residual and the uncertainty maps
deceptively benign.
\end{abstract}

\section*{Summary of findings}

\begin{enumerate}
\item \textbf{Two invariants per scan (exact).} Each pathway amplitude is
  $a_k=v\,G_k(u,E_2,\dw)$ with $u=(E_1-\cos\alpha)/(1-E_1\cos\alpha)$ and
  $v=\sin\alpha\,(1-E_1)/(1-E_1\cos\alpha)$ (\cref{thm:uv}). One scan identifies only
  $(u,\Mzero v)$: its Fisher information about $(\Mzero,\Bone,\Tone)$ is singular at every
  flip angle (\cref{cor:onescan}).
\item \textbf{Small flip angles are degenerate to leading order.} In the joint limit of small
  $\alpha$ and $\TR/\Tone$ the signal is invariant under
  $(\Mzero,\Bone,\Tone)\mapsto(\Mzero/\kappa,\kappa\Bone,\Tone/\kappa^2)$ (\cref{prop:scaling}).
  With $15^\circ/30^\circ$ the information matrix has condition number $5\times10^5$; the
  parameters remain identifiable in principle but not usefully at realistic SNR.
\item \textbf{The near-$360^\circ$ pulse breaks the degeneracy.} At $\Bone=1$, $330^\circ$ and
  $-30^\circ$ carry the same information (\cref{prop:period}); the benefit of $330^\circ$ is
  the eleven-fold faster response of its effective angle to $\Bone$. The condition number
  drops to 173 and the bound on $\log\Tone$ is $40.5\,\sigma/\Mzero$.
\item \textbf{Exact global aliases.} Magnitudes admit an exact alias with
  $\alpha_2=360^\circ+x$ inside the physiological range; complex data with a shared phase
  exclude it but admit one with $\alpha_2>540^\circ$, removed by the range assumption
  $\Bone<1.64$ (\cref{prop:equiv}).
\item \textbf{The published inverse is inefficient; ML is efficient.} The analytic inverse
  is exact without noise but its $\Tone$ standard deviation is 3.2--7.1 times the bound, with
  large bias and undefined estimates at moderate SNR. Joint ML attains the bound (ratio
  0.99--1.02) where its assumptions hold.
\item \textbf{Magnitudes need one bit of phase.} Magnitudes carry essentially all the local
  information, but their likelihood has two equal maxima for every noise realisation; the
  FID-sign test resolves it, after which magnitude least squares matches complex ML.
\item \textbf{Exact Jacobians make it fast.} Implicit differentiation of the steady state
  gives exact Jacobians at the cost of about one model evaluation; fits are 10--12 times
  faster in double precision and up to 34 times in single precision.
\item \textbf{Phantom.} On a $256\times256$ digital phantom (\Voxels{} brain voxels) the
  estimators behave as the voxel-wise analysis predicts, the published low-resolution design
  costs precision in $\Tone$ that a two-stage ML pipeline recovers in part, and the per-voxel
  uncertainty maps are calibrated (\cref{sec:phantom}).
\item \textbf{Partial volume.} On BrainWeb, WM/GM mixtures are harmless (second-order model
  error), but voxels containing CSF are biased by far more than the noise, with uncertainty
  intervals that cover the truth in as few as $0.5\%$ of voxels and residuals that a $\chi^2$ test
  flags in only $3$--$15\%$ at SNR 1000; the analytic inverse additionally jumps to a wrong flip-angle
  root in up to $38\%$ of them (\cref{sec:pv}).
\end{enumerate}

\section*{How to read this note}

\Cref{sec:notation} fixes notation. \Cref{sec:model} derives the forward model and proves
its structural properties. \Cref{sec:paper} presents and verifies the published inverse.
\Cref{sec:information} analyses what the data can determine, locally (CRLB, conditioning)
and globally (aliases). \Cref{sec:ml,sec:magnitude} develop and evaluate maximum-likelihood
estimation for complex data and least squares for magnitude data. \Cref{sec:jacobian}
derives exact Jacobians and their consequences. \Cref{sec:phantom} demonstrates everything
on a digital phantom, and \cref{sec:pv} measures the effect of partial volume on the
BrainWeb phantom. \Cref{sec:discussion} summarises which limitations belong to the
forward model and which to an inverse, and lists future work (model mismatch, Rician ML,
hierarchical estimation, learned estimators). Longer proofs are in \cref{app:proofs},
algorithms in \cref{app:algorithms}, and reproduction instructions in \cref{app:repro}.
